\section{Three rows}\label{sec:three}

On three rows the row bound of Theorem~\ref{thm:rowbound} is attained, and we count the
maximum blackouts as words. Throughout, $m\ge3$, the rows are $1,2,3$ from top to bottom, and the
tilted rectangles are the diamonds $D(i,1)=\{(i,2),(i+1,1),(i+2,2),(i+1,3)\}$ for
$1\le i\le m-2$ (Lemma~\ref{lem:tilted-narrow}).

\begin{theorem}[three rows]\label{thm:threevalue}
For every $m\ge3$, the maximum blackout of the $3\times m$ grid has $m+2$ points, that is,
$f(3,m)=m+2$.
\end{theorem}

\begin{proof}
The upper bound is the case $h=3$ of Theorem~\ref{thm:rowbound}, since $q_3=2$. The blackout
of row $1$ and column $1$ has $m+2$ points. It is the boundary cross of Lemma~\ref{lem:cross}
for $h=3$, so it is valid.
\end{proof}

As said in Section~\ref{sec:intro}, the trureturing project~\cite{trureturing2026} proved this
value earlier, by machine.

From now on every blackout has $m+2$ points. We name the columns by letters, $a=\{1\}$,
$b=\{2\}$ and $c=\{3\}$ for the singletons, $A=\{1,2\}$, $B=\{1,3\}$ and $C=\{2,3\}$ for the
doubles, and $X=\{1,2,3\}$ for the triple. The letters $a$ and $c$ are \emph{outer}. A blackout
with no empty column is the same thing as its \emph{word} $B_1B_2\cdots B_m$. Here $F$ is the
pair set of Section~\ref{sec:conv}, and we name each pair by its double letter.

\begin{proposition}[the two types]\label{prop:threeclass}
Let $m\ge3$ and let $S$ be a blackout of the $3\times m$ grid with $m+2$ points.
\begin{enumerate}[(a)]
\item If $S$ is valid, then its word has either exactly one letter $X$ and singletons elsewhere
(type T), or exactly two different doubles and singletons elsewhere (type P).
\item Conversely, for every word of type T or P the graph $G_S$ is a tree, so no two
axis-aligned rectangles have the same presentation.
\item There are exactly $m^23^{m-1}$ words of type T or P.
\end{enumerate}
\end{proposition}

\begin{proof}
(a) By Lemma~\ref{lem:H}, $|S|=t+e(H)$, where $t\le m$ counts the nonempty columns and the row
graph $H$ has at most two edges (Lemma~\ref{lem:mantel}). Both bounds are attained, so every
column is nonempty and $\sum_i(|B_i|-1)=2$. Two equal doubles form a four-cycle in $G_S$, which
Lemma~\ref{lem:girth} forbids.

(b) The column of $X$, or the two doubles, connect the three rows, and every other column hangs
on one row. So $G_S$ is a tree with $m+3$ vertices, and Lemma~\ref{lem:girth} applies.

(c) Type T gives $m3^{m-1}$ words, and type P gives $\binom m2\cdot6\cdot3^{m-2}=m(m-1)3^{m-1}$.
\end{proof}

We call the words of type T or P \emph{candidates}. They fall into the three \emph{classes}
$\{A,C\}$, $\{A,B\}$ and $\{B,C\}$ of type P and the triple class $X$. The left-to-right order of
the doubles splits a class of type P into two \emph{ordered classes}, such as $AC$ and $CA$. For
an ordered class such as $AC$, let $AC_m$ be its set of valid words of width $m$, and likewise
for the others.

\begin{lemma}[diamond test]\label{lem:diamond}
Let $S$ be a candidate with pair set $F$. Then $S$ is valid if and only if every diamond
$D(i,1)$ passes the test below, where $P=D(i,1)\setminus S$ is its presentation.
\begin{enumerate}[(a)]
\item If $|P|\ge3$ or $P$ is empty, the diamond passes.
\item If $P=\{(i+1,1),(i+1,3)\}$, the diamond fails if and only if $B\in F$.
\item If $P=\{(i,2),(i+2,2)\}$, it fails if and only if $B_i\cap B_{i+2}$ meets $\{1,3\}$.
\item If $P=\{(x,r),(z,s)\}$ with $x\ne z$, $r\ne s$, it fails if and only if $(x,s),(z,r)\in S$.
\item If $P=\{(x,r)\}$, it fails if and only if $\{r,s\}\in F$ for some $s\in B_x$.
\end{enumerate}
\end{lemma}

\begin{proof}
The axis-aligned rectangles are told apart (Proposition~\ref{prop:threeclass}(b)), and three
points fix a rectangle (Lemma~\ref{lem:threecorners}). Tests (b) to (e) are the cases of Lemma~\ref{lem:rival} for a diamond. In (b) a rival needs
another column black in rows $1$ and $3$, that is, $B\in F$. In (e) it needs a second column
black in rows $r$ and $s$, that is, $\{r,s\}\in F$. An empty diamond has no axis-aligned rival, since $G_S$ has no four-cycle. Two diamonds $D(i,1)$
and $D(j,1)$ with $i<j$ share at most one point, so a common presentation has at most one point.
Then columns $i+1$ and $j+1$ are both black in rows $1$ and $3$, which no candidate allows.
\end{proof}

A \emph{window} is a string of three consecutive letters, at columns $i$, $i+1$ and $i+2$. It
\emph{passes} or \emph{fails} with the diamond $D(i,1)$. We call $(i,2)$ and $(i+2,2)$ the
\emph{tips} of the window. A tip is black exactly when its letter is $b$, $A$, $C$ or $X$.

\begin{lemma}[singleton dichotomy]\label{lem:dichotomy}
Let $S$ be a candidate and let $stu$ be a window of three singletons.
\begin{enumerate}[(a)]
\item If $B\notin F$, the window passes.
\item If $B\in F$, the window fails if and only if it is $bab$, $bbb$ or $bcb$.
\end{enumerate}
\end{lemma}

\begin{proof}
If $s=u=b$, the diamond shows the vertical pair (for $t=b$) or one point of column $i+1$ (for
$t$ outer), and tests (b) and (e) fail it exactly when $B\in F$. If exactly one tip is black,
the diamond shows three corners or an oblique pair, and test (d) for that pair needs $(i+1,2)$
black, which it is not. If no tip is black, three corners show.
\end{proof}

\begin{lemma}[six of nine in the class $\{A,C\}$]\label{lem:sixofnine}
Let $S$ be a candidate of class $\{A,C\}$, so $F=\{A,C\}$, and let $s,t$ be singletons.
\begin{enumerate}[(a)]
\item A window $sAt$ fails if and only if $st$ is $bb$, $bc$ or $cb$.
\item A window $sCt$ fails if and only if $st$ is $bb$, $ba$ or $ab$.
\item A window with one double, at a tip, passes.
\end{enumerate}
\end{lemma}

\begin{proof}
(a) Here $(i+1,1)$ and $(i+1,2)$ are black and $(i+1,3)$ is white. For $s=t=b$ the diamond shows
$(i+1,3)$ alone, and test (e) fails it through $2\in B_{i+1}$ and $C\in F$. For $s=b\ne t$ it
shows $(i+2,2)$ and $(i+1,3)$, and test (d) asks whether $(i+2,3)$ is black, that is, whether
$t=c$. The case $t=b\ne s$ is the mirror image, and with no black tip three corners show.

(b) This is (a) after the reflection in the middle row, which swaps $a$ with $c$ and $A$ with $C$.

(c) Let the window be $Ast$, so the tip $(i,2)$ is black. If $t=b$, the diamond shows the
vertical pair or one point of column $i+1$, and both pass by tests (b) and (e), since
$B\notin F$. If $t\ne b$, it shows three corners or an oblique pair whose test (d)
needs $(i+1,2)$ black, and that point is white. Reflection and reversal give the other windows.
\end{proof}

\begin{figure}[t]
\centering
\resizebox{0.85\textwidth}{!}{\begin{tikzpicture}[x=0.52cm,y=0.52cm,baseline=(current bounding box.north)]
  \draw[gridcol] (1,-1) grid (3,-3);
  \fill[black] (1,-1) circle (0.17);
  \filldraw[fill=white,draw=keptline,line width=0.6pt] (1,-2) circle (0.17);
  \filldraw[fill=white,draw=keptline,line width=0.6pt] (1,-3) circle (0.17);
  \fill[black] (2,-1) circle (0.17);
  \fill[black] (2,-2) circle (0.17);
  \filldraw[fill=white,draw=keptline,line width=0.6pt] (2,-3) circle (0.17);
  \fill[black] (3,-1) circle (0.17);
  \filldraw[fill=white,draw=keptline,line width=0.6pt] (3,-2) circle (0.17);
  \filldraw[fill=white,draw=keptline,line width=0.6pt] (3,-3) circle (0.17);
  \draw[draw=rectA,line width=1.2pt,line join=round] (2,-1) -- (3,-2) -- (2,-3) -- (1,-2) -- cycle;
  \node[capt,anchor=north] at (2,-3.4) {$aAa$};
  \node[lab] at (0.35,-1) {1};
  \node[lab] at (0.35,-2) {2};
  \node[lab] at (0.35,-3) {3};
  \node[lab] at (1,-0.35) {$i$};
  \node[lab] at (2,-0.35) {$i+1$};
  \node[lab] at (3,-0.35) {$i+2$};
  \draw[gridcol] (3.9,-1) grid (5.9,-3);
  \fill[black] (3.9,-1) circle (0.17);
  \filldraw[fill=white,draw=keptline,line width=0.6pt] (3.9,-2) circle (0.17);
  \filldraw[fill=white,draw=keptline,line width=0.6pt] (3.9,-3) circle (0.17);
  \fill[black] (4.9,-1) circle (0.17);
  \fill[black] (4.9,-2) circle (0.17);
  \filldraw[fill=white,draw=keptline,line width=0.6pt] (4.9,-3) circle (0.17);
  \filldraw[fill=white,draw=keptline,line width=0.6pt] (5.9,-1) circle (0.17);
  \fill[black] (5.9,-2) circle (0.17);
  \filldraw[fill=white,draw=keptline,line width=0.6pt] (5.9,-3) circle (0.17);
  \draw[draw=rectA,line width=1.2pt,line join=round] (4.9,-1) -- (5.9,-2) -- (4.9,-3) -- (3.9,-2) -- cycle;
  \node[capt,anchor=north] at (4.9,-3.4) {$aAb$};
  \draw[gridcol] (6.8,-1) grid (8.8,-3);
  \fill[black] (6.8,-1) circle (0.17);
  \filldraw[fill=white,draw=keptline,line width=0.6pt] (6.8,-2) circle (0.17);
  \filldraw[fill=white,draw=keptline,line width=0.6pt] (6.8,-3) circle (0.17);
  \fill[black] (7.8,-1) circle (0.17);
  \fill[black] (7.8,-2) circle (0.17);
  \filldraw[fill=white,draw=keptline,line width=0.6pt] (7.8,-3) circle (0.17);
  \filldraw[fill=white,draw=keptline,line width=0.6pt] (8.8,-1) circle (0.17);
  \filldraw[fill=white,draw=keptline,line width=0.6pt] (8.8,-2) circle (0.17);
  \fill[black] (8.8,-3) circle (0.17);
  \draw[draw=rectA,line width=1.2pt,line join=round] (7.8,-1) -- (8.8,-2) -- (7.8,-3) -- (6.8,-2) -- cycle;
  \node[capt,anchor=north] at (7.8,-3.4) {$aAc$};
  \draw[gridcol] (9.7,-1) grid (11.7,-3);
  \filldraw[fill=white,draw=keptline,line width=0.6pt] (9.7,-1) circle (0.17);
  \fill[black] (9.7,-2) circle (0.17);
  \filldraw[fill=white,draw=keptline,line width=0.6pt] (9.7,-3) circle (0.17);
  \fill[black] (10.7,-1) circle (0.17);
  \fill[black] (10.7,-2) circle (0.17);
  \filldraw[fill=white,draw=keptline,line width=0.6pt] (10.7,-3) circle (0.17);
  \fill[black] (11.7,-1) circle (0.17);
  \filldraw[fill=white,draw=keptline,line width=0.6pt] (11.7,-2) circle (0.17);
  \filldraw[fill=white,draw=keptline,line width=0.6pt] (11.7,-3) circle (0.17);
  \draw[draw=rectA,line width=1.2pt,line join=round] (10.7,-1) -- (11.7,-2) -- (10.7,-3) -- (9.7,-2) -- cycle;
  \node[capt,anchor=north] at (10.7,-3.4) {$bAa$};
  \draw[gridcol] (12.6,-1) grid (14.6,-3);
  \filldraw[fill=white,draw=keptline,line width=0.6pt] (12.6,-1) circle (0.17);
  \fill[black] (12.6,-2) circle (0.17);
  \filldraw[fill=white,draw=keptline,line width=0.6pt] (12.6,-3) circle (0.17);
  \fill[black] (13.6,-1) circle (0.17);
  \fill[black] (13.6,-2) circle (0.17);
  \filldraw[fill=white,draw=keptline,line width=0.6pt] (13.6,-3) circle (0.17);
  \filldraw[fill=white,draw=keptline,line width=0.6pt] (14.6,-1) circle (0.17);
  \fill[black] (14.6,-2) circle (0.17);
  \filldraw[fill=white,draw=keptline,line width=0.6pt] (14.6,-3) circle (0.17);
  \draw[draw=rectB,line width=1.2pt,line join=round,dash pattern=on 3pt off 2pt] (13.6,-1) -- (14.6,-2) -- (13.6,-3) -- (12.6,-2) -- cycle;
  \node[capt,anchor=north] at (13.6,-3.4) {$bAb$};
  \draw[gridcol] (15.5,-1) grid (17.5,-3);
  \filldraw[fill=white,draw=keptline,line width=0.6pt] (15.5,-1) circle (0.17);
  \fill[black] (15.5,-2) circle (0.17);
  \filldraw[fill=white,draw=keptline,line width=0.6pt] (15.5,-3) circle (0.17);
  \fill[black] (16.5,-1) circle (0.17);
  \fill[black] (16.5,-2) circle (0.17);
  \filldraw[fill=white,draw=keptline,line width=0.6pt] (16.5,-3) circle (0.17);
  \filldraw[fill=white,draw=keptline,line width=0.6pt] (17.5,-1) circle (0.17);
  \filldraw[fill=white,draw=keptline,line width=0.6pt] (17.5,-2) circle (0.17);
  \fill[black] (17.5,-3) circle (0.17);
  \draw[draw=rectB,line width=1.2pt,line join=round,dash pattern=on 3pt off 2pt] (16.5,-1) -- (17.5,-2) -- (16.5,-3) -- (15.5,-2) -- cycle;
  \node[capt,anchor=north] at (16.5,-3.4) {$bAc$};
  \draw[gridcol] (18.4,-1) grid (20.4,-3);
  \filldraw[fill=white,draw=keptline,line width=0.6pt] (18.4,-1) circle (0.17);
  \filldraw[fill=white,draw=keptline,line width=0.6pt] (18.4,-2) circle (0.17);
  \fill[black] (18.4,-3) circle (0.17);
  \fill[black] (19.4,-1) circle (0.17);
  \fill[black] (19.4,-2) circle (0.17);
  \filldraw[fill=white,draw=keptline,line width=0.6pt] (19.4,-3) circle (0.17);
  \fill[black] (20.4,-1) circle (0.17);
  \filldraw[fill=white,draw=keptline,line width=0.6pt] (20.4,-2) circle (0.17);
  \filldraw[fill=white,draw=keptline,line width=0.6pt] (20.4,-3) circle (0.17);
  \draw[draw=rectA,line width=1.2pt,line join=round] (19.4,-1) -- (20.4,-2) -- (19.4,-3) -- (18.4,-2) -- cycle;
  \node[capt,anchor=north] at (19.4,-3.4) {$cAa$};
  \draw[gridcol] (21.3,-1) grid (23.3,-3);
  \filldraw[fill=white,draw=keptline,line width=0.6pt] (21.3,-1) circle (0.17);
  \filldraw[fill=white,draw=keptline,line width=0.6pt] (21.3,-2) circle (0.17);
  \fill[black] (21.3,-3) circle (0.17);
  \fill[black] (22.3,-1) circle (0.17);
  \fill[black] (22.3,-2) circle (0.17);
  \filldraw[fill=white,draw=keptline,line width=0.6pt] (22.3,-3) circle (0.17);
  \filldraw[fill=white,draw=keptline,line width=0.6pt] (23.3,-1) circle (0.17);
  \fill[black] (23.3,-2) circle (0.17);
  \filldraw[fill=white,draw=keptline,line width=0.6pt] (23.3,-3) circle (0.17);
  \draw[draw=rectB,line width=1.2pt,line join=round,dash pattern=on 3pt off 2pt] (22.3,-1) -- (23.3,-2) -- (22.3,-3) -- (21.3,-2) -- cycle;
  \node[capt,anchor=north] at (22.3,-3.4) {$cAb$};
  \draw[gridcol] (24.2,-1) grid (26.2,-3);
  \filldraw[fill=white,draw=keptline,line width=0.6pt] (24.2,-1) circle (0.17);
  \filldraw[fill=white,draw=keptline,line width=0.6pt] (24.2,-2) circle (0.17);
  \fill[black] (24.2,-3) circle (0.17);
  \fill[black] (25.2,-1) circle (0.17);
  \fill[black] (25.2,-2) circle (0.17);
  \filldraw[fill=white,draw=keptline,line width=0.6pt] (25.2,-3) circle (0.17);
  \filldraw[fill=white,draw=keptline,line width=0.6pt] (26.2,-1) circle (0.17);
  \filldraw[fill=white,draw=keptline,line width=0.6pt] (26.2,-2) circle (0.17);
  \fill[black] (26.2,-3) circle (0.17);
  \draw[draw=rectA,line width=1.2pt,line join=round] (25.2,-1) -- (26.2,-2) -- (25.2,-3) -- (24.2,-2) -- cycle;
  \node[capt,anchor=north] at (25.2,-3.4) {$cAc$};
\end{tikzpicture}
}
\caption{The nine windows $sAt$ of class $\{A,C\}$ with the diamond $D(i,1)$. Black dots are
blacked out. The diamond is solid where the window passes and dashed where it fails.}
\label{fig:sixofnine}
\end{figure}

We first give a heuristic, and Corollary~\ref{cor:threeasymp} makes it precise. By

Lemma~\ref{lem:sixofnine}, a double of the class $\{A,C\}$ passes for six of its nine pairs of
neighbors (Figure~\ref{fig:sixofnine}). This rule holds only in this class (compare
Lemma~\ref{lem:threewindows}). Suppose the two doubles avoid the end columns and are at least
three columns apart. Then only the two windows centered at them can fail, and each passes for
six of nine neighbor choices, independently. All but $O(m)$ positions of the doubles are of
this kind. The class holds the fraction $(m-1)/(3m)$ of the candidates. So about
$\frac13\cdot\bigl(\frac69\bigr)^2=\frac4{27}$ of the candidates are valid.

By Lemma~\ref{lem:diamond}, for $m\ge3$ the set $AC_m$ consists of the words $UAVCW$ of width
$m$, with $U$, $V$ and $W$ strings of singletons, whose windows all pass the test with
$F=\{A,C\}$. We use the same description for $m=2$ and put $G_{AC}(z)=\sum_{m\ge2}|AC_m|z^m$.
We define $G_{AB}$ and $G_X$ in the same way from the words $UAVBW$ and $UXW$, with $F=\{A,B\}$
and $F=\{A,B,C\}$. Put $R=1-3z$. The strings of $n\ge1$ singletons with a prescribed letter
next to a double have the generating function $z/R$.

\begin{proposition}[the class $\{A,C\}$]\label{prop:threeGAC}
The generating function of the ordered class $AC$ is
\begin{equation}\label{eq:GAC}
G_{AC}(z)=\frac{z^2-4z^3+10z^4-8z^5-3z^6}{(1-3z)^3},
\end{equation}
so $|AC_2|=1$, $|AC_3|=5$ and
\begin{equation}\label{eq:GACcoef}
|AC_m|=2(m-1)(3m+2)3^{m-5}\qquad(m\ge4).
\end{equation}
\end{proposition}

\begin{proof}
\emph{Step 1. The windows with both doubles.} In a window $tAC$ the points $(i+1,1)$ and
$(i+2,2)$ are black and $(i+1,3)$ is white. For $t=b$ the diamond shows $(i+1,3)$ alone, and
test (e) fails it. For $t\ne b$ it shows $(i,2)$ and $(i+1,3)$, and test (d) fails it exactly
when $t=c$. By the half turn $ACt$ passes only for $t=c$, and $AtC$ passes by tests (b) and (e).

\emph{Step 2. Adjacent doubles.} In $UACW$ every window passes except possibly the two of
Step~1 (Lemmas~\ref{lem:dichotomy}(a) and~\ref{lem:sixofnine}(c)). So $U$ ends in $a$ and $W$
begins with $c$, if nonempty, and this case contributes $z^2(1+z/R)^2$.

\emph{Step 3. One singleton between.} In $UAtCW$ the window $AtC$ passes. By
Lemma~\ref{lem:sixofnine}(a), the last letter of $U$ has $p_t=3,1,2$ allowed values for
$t=a,b,c$, and by (b) the first letter of $W$ has $q_t=2,1,3$. So this case contributes
$z^3\sum_t(1+p_tz/R)(1+q_tz/R)=z^3(3+12z/R+13z^2/R^2)$, since $\sum p_t=\sum q_t=6$ and
$\sum p_tq_t=13$.

\emph{Step 4. At least two singletons between.} In $UAv_1\cdots v_kCW$ with $k\ge2$ only
$uAv_1$ and $v_kCw$ can fail. If $U$ is empty, $v_1$ is free, and otherwise six of the nine
pairs $(u,v_1)$ pass. So each end contributes $3+6z/R$, the free middle letters give $1/R$, and
this case contributes $z^4(3+6z/R)^2/R=9z^4(1-z)^2/R^3$.

\emph{Step 5. The coefficients.} The three cases sum to~\eqref{eq:GAC}. The coefficient of $z^m$ in $z^kR^{-3}$ is
$\binom{m-k+2}23^{m-k}$ for $m\ge k$ and $0$ otherwise. For $m\ge4$ the five terms with
$k=2,\dots,6$ sum to the polynomial $2(m-1)(3m+2)3^{m-5}$, and at $m=2,3$ they give $1$ and $5$.
\end{proof}

\begin{theorem}[the number of three-row optima]\label{thm:threecount}
For $m\ge3$ let $E_m$ be the number of maximum blackouts of the $3\times m$ grid whose word is
outside $AC$ and $CA$. Then for every $m\ge4$,
\begin{equation}\label{eq:threecount}
a_3(m)=4(m-1)(3m+2)3^{m-5}+E_m.
\end{equation}
\end{theorem}

\begin{proof}
The reversal of the columns maps $AC$ onto $CA$, $AB$ onto $BA$ and $BC$ onto $CB$. The
reflection in the middle row maps $AB$ onto $CB$. Both map optima to optima, so
\begin{equation}\label{eq:classsum}
a_3(m)=2|AC_m|+4|AB_m|+|X_m|=2|AC_m|+E_m\qquad(m\ge3),
\end{equation}
and~\eqref{eq:GACcoef} gives~\eqref{eq:threecount}.
\end{proof}

Note that $E_m$ is also the number of optima with a column black in rows $1$ and $3$, as in
Theorem~B, since such a column is $B$ or $X$. At $m=3$ the right side of~\eqref{eq:GACcoef} is
$44/9$, which is not an integer, so~\eqref{eq:threecount} holds only from $m=4$ on. Directly,
$a_3(3)=53=2\cdot5+4\cdot7+15$, with $|AB_3|=7$ and $|X_3|=15$ from
Propositions~\ref{prop:GAB} and~\ref{prop:GX}.

\subsection{The generating function}

Put $L=1-2z-2z^2$, $N=1+2z^2$ and $D=LN$.

\begin{theorem}[three-row generating function, parts (b) and (c) computer-assisted]\label{thm:threegf}
\leavevmode
\begin{enumerate}[(a)]
\item The series $\sum_{m\ge3}a_3(m)z^m$ equals $2G_{AC}+4G_{AB}+G_X-z-12z^2$. Here $G_{AC}$
is~\eqref{eq:GAC}, and $G_{AB}$ and $G_X$ are the rational functions of
Propositions~\ref{prop:GAB} and~\ref{prop:GX}, with denominators dividing $D^3$ and $D^2$.
\item The reduced denominator of the series is exactly
\[
Q(z)=(1-3z)^3(1-2z-2z^2)^3(1+2z^2)^3.
\]
\item With
\begin{align*}
P(z)={}&z-3z^2-34z^3+228z^4-641z^5+1307z^6-2186z^7\\
&+2056z^8-236z^9-1060z^{10}+2584z^{11}-4264z^{12}\\
&-1440z^{13}+704z^{14}-1024z^{15}+2080z^{16}+2176z^{17}+384z^{18},
\end{align*}
we have
\begin{equation}\label{eq:PQ}
\sum_{m\ge3}a_3(m)z^m=\frac{P(z)}{Q(z)}-z-12z^2.
\end{equation}
\end{enumerate}
\end{theorem}

\begin{proof}
(a) Use~\eqref{eq:classsum}. The coefficients $1$ and $12$ of $z$ and $z^2$ count the word $X$,
the six words $sX$ and $Xs$, and the six words of two different doubles.

(b) This is Theorem~\ref{thm:denominator} in Appendix~\ref{app:threegf}. The orders at the
roots of $L$ and $N$ rest on the computer-assisted form~\eqref{eq:PB} of $G_{AB}$.

(c) Clearing denominators in (a) gives $P$, which is coprime to $Q$. The script
\texttt{three\_row\_closed\_form.py} does this algebra exactly in SymPy. The script
\texttt{three\_row\_closed\_form\_check.py} checks~\eqref{eq:PQ} against an independent transfer
matrix with $350$ states, in exact integer arithmetic through width $100$.
\end{proof}

By~\eqref{eq:PQ}, $Q\sum_{m\ge3}a_3(m)z^m=P-(z+12z^2)Q$ has degree $18$ and leading
coefficient $384$. So the recurrence of order $15$ with the coefficients of $Q$ holds for
every $m\ge19$ and fails at $m=18$.

\begin{corollary}[asymptotics]\label{cor:threeasymp}
Let $\rho=1+\sqrt3$. As $m\to\infty$,
\begin{equation}\label{eq:asymp1}
a_3(m)=\frac4{81}m^23^m-\frac4{243}(m+2)3^m+O(m^2\rho^m).
\end{equation}
So the valid fraction $a_3(m)/(m^23^{m-1})$ of the candidates is
$\frac4{27}-\frac4{81m}-\frac8{81m^2}+O\bigl((\rho/3)^m\bigr)$.
\end{corollary}

\begin{proof}
The main term of~\eqref{eq:threecount} is
$4(3m^2-m-2)3^{m-5}=\frac4{81}m^23^m-\frac4{243}(m+2)3^m$. By Theorem~\ref{thm:threegf}(a)
and~\eqref{eq:classsum}, $E_m$ is the coefficient of $z^m$ in $4G_{AB}+G_X$, whose poles are roots of $D$ of order at
most three. The root of least modulus is $1/\rho$, so $E_m=O(m^2\rho^m)$.
\end{proof}

\begin{theorem}[three-row closed form, computer-assisted exact algebra]\label{thm:threeclosed}
Let $\sigma=1-\sqrt3$. For every $m\ge3$,
\begin{equation}\label{eq:Eclosed}
E_m=p_+(m)\rho^m+p_-(m)\sigma^m+T_m,
\end{equation}
where
\[
p_\pm(m)=\frac{(19\pm11\sqrt3)m^2+(61\pm41\sqrt3)m+(104\pm28\sqrt3)}{432},
\]
and $T_m=(-2)^{\lfloor m/2\rfloor}\,\tau(m)$ with $\tau(m)=(2m^2+41m-104)/216$ for even $m$ and
$\tau(m)=(7m^2-32m-175)/432$ for odd $m$. Hence $E_m\sim\frac{19+11\sqrt3}{432}m^2\rho^m$, and $E_m>0$ for every $m\ge4$.
\end{theorem}

\begin{proof}
Let $c_m$ be the coefficient of $z^m$ in $4G_{AB}+G_X$, a rational function with denominator
$L^3N^3$. By Theorem~\ref{thm:threegf}(a), $E_m=c_m$ for $m\ge3$. The roots $1/\rho$, $1/\sigma$
and $\pm1/(i\sqrt2)$ of $L^3N^3$ are triple, so beyond a polynomial part of degree one the
$c_m$ have the form~\eqref{eq:Eclosed}. The script \texttt{three\_row\_closed\_form.py} solves for
the twelve rational coefficients in SymPy. Both $c_m$ and the right side of~\eqref{eq:Eclosed}
satisfy the order-$12$ recurrence of $L^3N^3$ from $m=14$ on. So agreement at $m=2,\dots,13$
proves that they are equal for every $m\ge2$. The script \texttt{three\_row\_closed\_form\_check.py}
confirms it over $\mathbb Q(\sqrt3)$ for $m\le100$. The asymptotic follows since $|\sigma|<1$
and $|T_m|=O(m^22^{m/2})$. Finally $E_m>0$, since the cross $Xaa\cdots a$ of
Theorem~\ref{thm:threevalue} lies in the triple class.
\end{proof}


The values at widths $3$ to $12$ are $53$, $224$, $892$, $3420$, $12704$, $45864$, $161992$,
$562912$, $1930312$ and $6545496$. They start row $3$ of \seqnum{A399580}. The ILP program and
the brute-force check of Section~\ref{sec:verify} find them through width $9$. A search over
every blackout of size $m+2$ gives the class counts at widths $3$ to $10$
(\texttt{three\_row\_definition\_check.json}), and they agree with the series.

\begin{remark}[optima up to symmetry, computer-assisted]\label{rem:threeorbits}
Up to reflection and reversal, \texttt{three\_row\_orbits.py} finds $15$, $58$, $227$, $864$,
$3190$, $11502$, $40549$ and $140863$ optima at widths $3$ to $10$. For even $m$ from $4$ to
$14$ only the half turn fixes optima, exactly $2m3^{m/2-2}$ of them. We do not prove this here.

\end{remark}
